\documentclass[11pt]{article}
\usepackage[utf8]{inputenc}
\usepackage[T1]{fontenc}
\usepackage[spanish]{babel}
\usepackage{amsmath,amssymb,amsthm}
\usepackage{geometry}
\geometry{margin=1in}

\title{Resumen de la via cuartica auto-rec\'iproca hacia la funci\'on xi}
\author{}
\date{}

\begin{document}
\maketitle

\section*{1. Objetivo}

Queremos usar una familia cuartica basada solamente en
\[
t^4+t^{-4}
\]
para construir una regularizacion de la funcion xi de Riemann sin perder la
linea critica
\[
s=\frac12+iz.
\]
La cuestion no es afirmar que esto demuestra RH, sino identificar si la
estructura cuartica produce algo mas que la simetria funcional conocida.

\section*{2. Bloque cuartico basico}

El bloque crudo es
\[
I_{m,a}(z)
=
\int_0^\infty
t^{2iz}
e^{-\pi(m^2t^4+a^2t^{-4})}\,dt,
\qquad m\ge1,\ a>0.
\]
Con el cambio \(y=t^4\),
\[
I_{m,a}(z)
=
\frac12
\left(\frac{a}{m}\right)^{\frac14+\frac{iz}{2}}
K_{\frac14+\frac{iz}{2}}(2\pi ma).
\]
Aqui aparece de forma natural el orden
\[
\nu=\frac14+\frac{iz}{2}.
\]

En el limite formal \(a\downarrow0\),
\[
I_{m,0}(z)
=
\frac14
\pi^{-\nu}\Gamma(\nu)m^{-2\nu},
\qquad
\nu=\frac14+\frac{iz}{2}.
\]
Como
\[
2\nu=\frac12+iz=s,
\]
se obtiene
\[
2\sum_{m\ge1}I_{m,0}(z)
=
\frac12
\pi^{-s/2}\Gamma(s/2)\zeta(s).
\]
Por tanto,
\[
\xi(s)
=
\frac12s(s-1)\pi^{-s/2}\Gamma(s/2)\zeta(s)
=
s(s-1)\left(2\sum_{m\ge1}I_{m,0}(z)\right),
\]
con
\[
s=\frac12+iz,
\qquad
s(s-1)=-(z^2+\tfrac14).
\]

\section*{3. Por que el exponente correcto es \(t^{2iz}\)}

El exponente debe mantener
\[
s=\frac12+iz.
\]
Si se elige mal, se llega a lineas laterales como
\[
\zeta\left(\frac14+\frac{iz}{2}\right)
\quad\text{o}\quad
\zeta\left(\frac32+2iz\right).
\]
La linea \(\Re(s)=3/2\) es cero-libre pero no sirve para RH. La linea
\(\Re(s)=1/4\) no es segura: excluir ceros alli seria ya contenido fuerte
relacionado con RH. Por tanto hay que conservar exactamente
\[
s=\frac12+iz.
\]

\section*{4. Factor critico de auto-reciprocidad}

La transformada
\[
F(z)=\int_0^\infty t^{2iz}V(t)\,dt
\]
sera par, \(F(z)=F(-z)\), si el nucleo cumple
\[
V(t)=t^{-2}V(1/t).
\]
El factor \(t^{-2}\) no es decorativo: es el jacobiano de \(t\mapsto1/t\) y
es el factor que protege la geometria critica.

Por tanto, el bloque auto-reciproco critico es
\[
\widetilde V_{m,a}(t)
=
\frac12\left[
e^{-\pi(m^2t^4+a^2t^{-4})}
+
t^{-2}e^{-\pi(m^2t^{-4}+a^2t^4)}
\right].
\]
Entonces
\[
\widetilde V_{m,a}(t)=t^{-2}\widetilde V_{m,a}(1/t),
\]
y
\[
\widetilde I_{m,a}(z)
=
\int_0^\infty t^{2iz}\widetilde V_{m,a}(t)\,dt
=
\frac12\left(I_{m,a}(z)+I_{m,a}(-z)\right).
\]
En forma cerrada,
\[
\widetilde I_{m,a}(z)
=
\frac14\left[
\left(\frac{a}{m}\right)^{\frac14+\frac{iz}{2}}
K_{\frac14+\frac{iz}{2}}(2\pi ma)
+
\left(\frac{a}{m}\right)^{\frac14-\frac{iz}{2}}
K_{\frac14-\frac{iz}{2}}(2\pi ma)
\right].
\]

\section*{5. Funcion regularizada centrada}

La version centrada natural es
\[
\widetilde F_a(z)=2\sum_{m\ge1}\widetilde I_{m,a}(z),
\]
y
\[
\widetilde X_a(z)
=
-(z^2+\tfrac14)\widetilde F_a(z).
\]
Formalmente,
\[
\widetilde X_a(z)\longrightarrow \xi\left(\frac12+iz\right)
\]
cuando \(a\downarrow0\), despues de tratar los terminos polares o
elementales explicitos que aparecen en la regularizacion.

\section*{6. Identidad exacta de Mobius}

La identidad de eliminacion exacta es
\[
\sum_{d\ge1}\mu(d)\bigl(\vartheta(d^2x)-1\bigr)
=2e^{-\pi x}.
\]
Aplicada con \(x=a^2t^{-4}\),
\[
\boxed{
\sum_{d\ge1}\mu(d)
\bigl(\vartheta(d^2a^2t^{-4})-1\bigr)
=
2e^{-\pi a^2t^{-4}}.
}
\]
Esto elimina exactamente todos los modos del segundo theta salvo el modo
unitario. Esta es una parte no trivial de la arquitectura: no es simplemente
poner un filtro externo, sino una proyeccion aritmetica exacta dentro del
nucleo cuartico.

\section*{7. Que es estandar y que no}

Es estandar:
\[
\text{``hacer un nucleo auto-reciproco para obtener una transformada par''.}
\]
Tambien es estandar que el jacobiano correcto bajo \(t\mapsto1/t\) sea
\[
t^{-2}.
\]

No es tan estandar la combinacion simultanea:
\[
\text{Mobius exacto}
+
\text{regularizacion }t^4+t^{-4}
+
\text{auto-reciprocidad critica}
+
\text{limite hacia }\xi(1/2+iz).
\]
La pregunta real es si esta arquitectura produce una restriccion nueva sobre
ceros, no solo la simetria funcional.

\section*{8. Obstaculo principal: ceros de los bloques}

Para \(z\in\mathbb R\),
\[
\widetilde I_{m,a}(z)
=
\Re\left[
\frac12
\left(\frac{a}{m}\right)^{\frac14+\frac{iz}{2}}
K_{\frac14+\frac{iz}{2}}(2\pi ma)
\right].
\]
Si el orden fuera puramente imaginario, \(K_{i\tau}(x)\), hay teoria clasica
favorable. Pero aqui el orden esta desplazado:
\[
\nu=\frac14+\frac{iz}{2}.
\]
Ese desplazamiento real \(1/4\) mantiene la linea critica, pero rompe la
situacion mas favorable. No hay razon inmediata para que
\[
\Re\left[A^{1/4+iz/2}K_{1/4+iz/2}(x)\right]
\]
tenga solo ceros reales como funcion de \(z\). Demostrarlo seria muy dificil
y probablemente falso en general.

Incluso si cada bloque individual tuviera ceros reales, la suma
\[
\sum_{m\ge1}\widetilde I_{m,a}(z)
\]
no preservaria automaticamente esa propiedad. La suma de funciones con ceros
reales no tiene, en general, ceros reales.

Conclusion: no parece razonable intentar probar que
\[
\widetilde X_a
\]
tiene solo ceros reales para todo \(a>0\).

\section*{9. Flujo en \(a\)}

Para el bloque crudo,
\[
\partial_a I_{m,a}(z)
=
-2\pi a\, I_{m,a}(z+2i),
\]
salvo normalizaciones si se incluye el factor polinomico de completado.

Despues de simetrizar,
\[
\widetilde I_{m,a}(z)
=\frac12(I_{m,a}(z)+I_{m,a}(-z)),
\]
el flujo deja de ser escalar y se vuelve bilateral:
\[
z\mapsto z+2i,
\qquad
z\mapsto z-2i.
\]
Esto conserva la geometria critica, pero introduce un sistema par/impar o
bilateral. La dificultad no es ya mantener la linea \(1/2\); eso lo hace el
factor \(t^{-2}\). La dificultad es extraer de este sistema una condicion
nueva sobre los ceros.

\section*{10. Linea quarter}

La version quarter aparece al reescalar
\[
w=2s.
\]
Si
\[
w=\frac12+i\gamma,
\]
entonces
\[
s=\frac{w}{2}
=
\frac14+\frac{i\gamma}{2}.
\]
Esto no cambia la dificultad: solo desplaza la geometria de la linea
\(\Re(w)=1/2\) a
\[
\Re(s)=1/4.
\]
Lo peligroso es confundir esta linea quarter reescalada con la linea
\(\Re(s)=1/4\) de la zeta original. Excluir ceros de
\[
\zeta\left(\frac14+it\right)
\]
no es conocido y seria contenido fuerte relacionado con RH.

\section*{11. Diagnostico}

La modificacion correcta de \(V\) es:
\[
\boxed{
\widetilde V_{m,a}(t)
=
\frac12\left[
e^{-\pi(m^2t^4+a^2t^{-4})}
+
t^{-2}e^{-\pi(m^2t^{-4}+a^2t^4)}
\right].
}
\]
Esta eleccion:
\begin{itemize}
\item mantiene el centro critico \(s=1/2+iz\);
\item mantiene la estructura cuartica \(t^4+t^{-4}\);
\item permite usar la eliminacion exacta de Mobius;
\item produce una familia regularizada que converge formalmente a
  \(\xi(1/2+iz)\).
\end{itemize}

Pero por ahora no da una demostracion de RH. El obstaculo principal no es la
simetria, sino encontrar una restriccion adicional sobre ceros a partir de:
\[
\text{bloques }K_{1/4+iz/2},
\qquad
\text{flujo bilateral }z\mapsto z\pm2i,
\qquad
\text{proyeccion de Mobius}.
\]

\section*{12. Preguntas concretas para consultar}

\begin{enumerate}
\item \textbf{Bloques individuales.} Hay resultados conocidos sobre ceros
reales de
\[
\Re\left[A^{1/4+iz/2}K_{1/4+iz/2}(x)\right]?
\]
Si no, hay contraejemplos numericos?

\item \textbf{Suma en \(m\).} La suma
\[
\sum_{m\ge1}\widetilde I_{m,a}(z)
\]
pertenece a alguna clase de funciones con restricciones de ceros, o
probablemente no?

\item \textbf{Flujo bilateral.} La ecuacion en \(a\) con desplazamientos
\[
z\mapsto z\pm2i
\]
tiene alguna cantidad conservada, desigualdad de fase, o identidad tipo Turan
que no sea equivalente a RH?

\item \textbf{Terminos elementales.} Al simetrizar y usar theta, aparecen
correcciones elementales cuarticas/polares. Se pueden sustraer de forma
canonica para que el limite sea exactamente \(\xi(1/2+iz)\)?

\item \textbf{Originalidad.} La combinacion
\[
\text{Mobius exacto}+\text{cuartica auto-reciproca critica}
\]
esta ya estudiada en alguna forma conocida?
\end{enumerate}

\section*{13. Veredicto corto}

La via no se cae por la escala: el factor \(t^{-2}\) arregla exactamente el
centro critico. Tampoco es solo una multiplicacion externa de \(\xi\). La
parte prometedora es la identidad exacta de Mobius dentro del nucleo cuartico.

El muro visible es otro: no hay, de momento, un mecanismo claro que convierta
esa regularizacion en control de ceros. En particular, la estrategia
\[
\widetilde X_a \text{ tiene todos sus ceros reales para } a>0
\]
parece demasiado fuerte y probablemente falsa.

La ruta razonable es usar \(\widetilde X_a\) como regularizacion centrada de
\(\xi\), no como una familia RH intermedia, y buscar una identidad nueva en el
limite \(a\downarrow0\).

\end{document}
